\begin{textsample}{POS Dim 1 – 2020 – Score 32.00 – 2020/t004244.txt}  \label{ex:f1_pos_072}
Chris Anderson : Al, welcome.
So look, just six months ago -- it seems a lifetime ago, but it really was just six months ago -- climate seemed to be on \textbf{the} lips of every thinking person on \textbf{the} \textbf{planet}.
Recent events seem to have swept it all away from our attention.
How worried are you about that?
Al Gore : Well, first of all Chris, thank you so much for inviting me to have this conversation.
People are reacting differently to \textbf{the} climate crisis in \textbf{the} midst of these other great challenges that have taken over our awareness, appropriately.
One reason is something that you mentioned.
People get \textbf{the} fact that when scientists are warning us in ever more dire terms and setting their hair on fire, so to speak, it ' s best to listen to what they ' re saying, and I think that lesson has begun to sink in in a new way.
Another similarity, by \textbf{the} way, is that \textbf{the} climate crisis, like \textbf{the} COVID-19 pandemic, has revealed in a new way \textbf{the} shocking injustices and inequalities and disparities that affect communities of color and low-income communities.
There are differences. \textbf{The} climate crisis has effects that are not measured in years, as \textbf{the} pandemic is, but consequences that are measured in centuries and even longer.
And \textbf{the} other difference is that instead of depressing economic activity to deal with \textbf{the} climate crisis, as nations around \textbf{the} world have had to do with COVID-19, we have \textbf{the} opportunity to create tens of millions of new jobs.
That sounds like a political phrasing, but it ' s literally true.
For \textbf{the} last five years, \textbf{the} fastest-growing job in \textbf{the} US has been solar installer. \textbf{The} second-fastest has been wind \textbf{turbine} technician.
And \textbf{the}"Oxford Review of Economics,"just a few weeks ago, pointed \textbf{the} way to a very jobs-rich recovery if we emphasize \textbf{renewable} energy and \textbf{sustainability} technology.
So I think we are crossing a tipping point, and you need only look at \textbf{the} recovery plans that are being presented in nations around \textbf{the} world to see that they ' re very much focused on a green recovery.
CA : I mean, one obvious impact of \textbf{the} pandemic is that it ' s brought \textbf{the} world ' s economy to a shuddering halt, thereby reducing \textbf{greenhouse} gas \textbf{emissions}.
I mean, how big an effect has that been, and is it unambiguously good news?
AG : Well, it ' s a little bit of an illusion, Chris, and you need only look back to \textbf{the} Great Recession in 2008 and ' 09, when there was a one percent decline in \textbf{emissions}, but then in 2010, they came roaring back during \textbf{the} recovery with a four percent increase. \textbf{The} latest estimates are that \textbf{emissions} will go down by at least five percent during this induced coma, as \textbf{the} economist Paul Krugman perceptively described it, but whether it goes back \textbf{the} way it did after \textbf{the} Great Recession is in part up to us, and if these green recovery plans are actually implemented, and I know many countries are determined to implement them, then we need not repeat that pattern.
After all, this whole process is \textbf{occurring} during a period when \textbf{the} cost of \textbf{renewable} energy and electric vehicles, batteries and a range of other \textbf{sustainability} approaches are continuing to fall in price, and they ' re becoming much more competitive.
Just a quick reference to how fast this is : five years ago, electricity from solar and wind was cheaper than electricity from \textbf{fossil} fuels in only one percent of \textbf{the} world.
This year, it ' s cheaper in \textbf{two-thirds} of \textbf{the} world, and five years from now, it will be cheaper in virtually 100 percent of \textbf{the} world.
EVs will be cost-competitive within two years, and then will continue falling in price.
And so there are changes underway that could interrupt \textbf{the} pattern we saw after \textbf{the} Great Recession.
CA : \textbf{The} reason those pricing differentials happen in different parts of \textbf{the} world is obviously because there ' s different amounts of sunshine and wind there and different building costs and so forth.
AG : Well, yes, and government policies also account for a lot. \textbf{The} world is continuing to subsidize \textbf{fossil} fuels at a ridiculous amount, more so in many developing countries than in \textbf{the} US and developed countries, but it ' s subsidized here as well.
But everywhere in \textbf{the} world, wind and solar will be cheaper as a source of electricity than \textbf{fossil} fuels, within a few years.
CA : I think I ' ve heard it said that \textbf{the} fall in \textbf{emissions} caused by \textbf{the} pandemic isn ' t that much more than, actually, \textbf{the} fall that we will need every single year if we ' re to meet \textbf{emissions} targets.
Is that true, and, if so, doesn ' t that seem impossibly daunting?
AG : It does seem daunting, but first look at \textbf{the} number.
That number came from a study a little over a year ago released by \textbf{the} IPCC as to what it would take to keep \textbf{the} Earth ' s temperatures from increasing more than 1. 5 degrees Celsius.
And yes, \textbf{the} annual reductions would be significant, on \textbf{the} order of what we ' ve seen with \textbf{the} pandemic.
And yes, that does seem daunting.
However, we do have \textbf{the} opportunity to make some fairly dramatic changes, and \textbf{the} plan is not a mystery.
You start with \textbf{the} two sectors that are closest to an effective transition -- electricity generation, as I mentioned -- and last year, 2019, if you look at all of \textbf{the} new electricity generation built all around \textbf{the} world, 72 percent of it was from solar and wind.
And already, without \textbf{the} continuing subsidies for \textbf{fossil} fuels, we would see many more of these plants being shut down.
There are some new \textbf{fossil} plants being built, but many more are being shut down.
And where \textbf{transportation} is concerned, \textbf{the} second sector ready to go, in addition to \textbf{the} cheaper prices for EVs that I made reference to before, there are some 45 jurisdictions around \textbf{the} world -- national, regional and municipal -- where laws have been passed beginning a phaseout of internal combustion engines.
Even India said that by 2030, less than 10 years from now, it will be illegal to sell any new internal combustion engines in India.
There are many other examples.
So \textbf{the} past small reductions may not be an accurate guide to \textbf{the} kind we can achieve with serious national plans and a focused global effort.
CA : So help us understand just \textbf{the} big picture here, Al.
I think before \textbf{the} pandemic, \textbf{the} world was emitting about 55 gigatons of what they call"CO2 equivalent,"so that includes other \textbf{greenhouse} gases like methane dialed up to be \textbf{the} equivalent of CO2.
And am I right in saying that \textbf{the} IPCC, which is \textbf{the} global organization of scientists, is recommending that \textbf{the} only way to fix this crisis is to get that number from 55 to zero by 2050 at \textbf{the} very latest, and that even then, there ' s a chance that we will end up with temperature rises more like two degrees Celsius rather than 1. 5?
I mean, is that approximately \textbf{the} big picture of what \textbf{the} IPCC is recommending?
AG : That ' s correct. \textbf{The} global goal established in \textbf{the} Paris Conference is to get to net zero on a global basis by 2050, and many people quickly add that that really means a 45 to 50 percent reduction by 2030 to make that pathway to net zero feasible.
CA : And that kind of timeline is \textbf{the} kind of timeline where people couldn ' t even imagine it.
It ' s just hard to think of policy over 30 years.
So that ' s actually a very good shorthand, that humanity ' s task is to cut \textbf{emissions} in half by 2030, approximately speaking, which I think boils down to about a seven or eight percent reduction a year, something like that, if I ' m not wrong.
AG : Not quite.
Not quite that large but close, yes.
CA : So it is something like \textbf{the} effect that we ' ve experienced this year may be necessary.
This year, we ' ve done it by basically shutting down \textbf{the} economy.
You ' re talking about a way of doing it over \textbf{the} coming years that actually gives some economic growth and new jobs.
So talk more about that.
You ' ve referred to changing our energy sources, changing how we transport.
If we did those things, how much of \textbf{the} problem does that solve?
AG : Well, we can get to -- well, in addition to doing \textbf{the} two sectors that I mentioned, we also have to deal with manufacturing and all \textbf{the} use cases that require temperatures of a thousand degrees Celsius, and there are solutions there as well.
I ' ll come back and mention an exciting one that Germany has just embarked upon.
We also have to tackle regenerative \textbf{agriculture}.
There is \textbf{the} opportunity to sequester a great deal of \textbf{carbon} in topsoils around \textbf{the} world by changing \textbf{the} agricultural \textbf{techniques}.
There is a farmer-led movement to do that.
We need to also retrofit buildings.
We need to change our management of forests and \textbf{the} ocean.
But let me just mention two things briefly.
First of all, \textbf{the} high temperature use cases.
Angela Merkel, just 10 days ago, with \textbf{the} leadership of her minister Peter Altmaier, who is a good friend and a great public servant, have just embarked on a green hydrogen strategy to make hydrogen with zero marginal cost \textbf{renewable} energy.
And just a word on that, Chris : you ' ve heard about \textbf{the} intermittency of wind and solar -- solar doesn ' t produce electricity when \textbf{the} sun ' s not shining, and wind doesn ' t when \textbf{the} wind ' s not blowing -- but batteries are getting better, and these technologies are becoming much more efficient and powerful, so that for an increasing number of hours of each day, they ' re producing often way more electricity than can be used.
So what to do with it? \textbf{The} marginal cost for \textbf{the} next kilowatt-hour is zero.
So all of a sudden, \textbf{the} very energy-intensive process of cracking hydrogen from water becomes economically feasible, and it can be substituted for \textbf{coal} and gas, and that ' s already being done.
There ' s a Swedish company already making steel with green hydrogen, and, as I say, Germany has just embarked on a major new initiative to do that.
I think they ' re pointing \textbf{the} way for \textbf{the} rest of \textbf{the} world.
Now, where building retrofits are concerned, just a moment on this, because about 20 to 25 percent of \textbf{the} global warming pollution in \textbf{the} world and in \textbf{the} US comes from inefficient buildings that were constructed by companies and individuals who were trying to be competitive in \textbf{the} marketplace and keep their margins acceptably high and thereby skimping on insulation and \textbf{the} right windows and LEDs and \textbf{the} rest.
And yet \textbf{the} person or company that buys that building or leases that building, they want their monthly utility bills much lower.
So there are now ways to close that so-called agent-principal divide, \textbf{the} differing incentives for \textbf{the} builder and occupier, and we can retrofit buildings with a program that literally pays for itself over three to five years, and we could put tens of millions of people to work in jobs that by definition cannot be outsourced because they exist in every single community.
And we really ought to get serious about doing this, because we ' re going to need all those jobs to get sustainable prosperity in \textbf{the} aftermath of this pandemic.
CA : Just going back to \textbf{the} hydrogen economy that you referred to there, when some people hear that, they think,"Oh, are you talking about hydrogen-fueled cars?"And they ' ve heard that that probably won ' t be a winning strategy.
But you ' re thinking much more broadly than that, I think, that it ' s not just hydrogen as a kind of storage mechanism to act as a buffer for \textbf{renewable} energy, but also hydrogen could be essential for some of \textbf{the} other processes in \textbf{the} economy like making steel, making cement, that are fundamentally carbon-intensive processes right now but could be transformed if we had much cheaper sources of hydrogen.
Is that right?
AG : Yes, I was always skeptical about hydrogen, Chris, principally because it ' s been so expensive to make it, to"crack it out of water,"as they say.
But \textbf{the} game-changer has been \textbf{the} incredible abundance of solar and wind electricity in volumes and amounts that people didn ' t expect, and all of a sudden, it ' s cheap enough to use for these very energy-intensive processes like creating green hydrogen.
I ' m still a bit skeptical about using it in vehicles.
Toyota ' s been betting on that for 25 years and it hasn ' t really worked for them.
Never say never, maybe it will, but I think it ' s most useful for these high-temperature industrial processes, and we already have a pathway for decarbonizing \textbf{transportation} with electricity that ' s working extremely well.
Tesla ' s going to be soon \textbf{the} most valuable automobile company in \textbf{the} world, already in \textbf{the} US, and they ' re about to overtake Toyota.
There is now a semitruck company that ' s been stood up by Tesla and another that is going to be a hybrid with electricity and green hydrogen, so we ' ll see whether or not they can make it work in that application.
But I think electricity is preferable for cars and trucks.
CA : We ' re coming to some community questions in a minute.
Let me ask you, though, about nuclear.
Some environmentalists believe that nuclear, or maybe new generation nuclear power is an essential part of \textbf{the} equation if we ' re to get to a truly clean future, a clean energy future.
Are you still pretty skeptical on nuclear, Al?
AG : Well, \textbf{the} market ' s skeptical about it, Chris.
It ' s been a crushing disappointment for me and for so many.
I used to represent Oak Ridge, where nuclear energy began, and when I was a young congressman, I was a booster.
I was very enthusiastic about it.
But \textbf{the} cost overruns and \textbf{the} problems in building these plants have become so severe that utilities just don ' t have an appetite for them.
It ' s become \textbf{the} most expensive source of electricity.
Now, let me hasten to add that there are some older nuclear reactors that have more useful time that could be added onto their lifetimes.
And like a lot of environmentalists, I ' ve come to \textbf{the} view that if they can be determined to be safe, they should be allowed to continue operating for a time.
But where new nuclear power plants are concerned, here ' s a way to look at it.
If you are -- you ' ve been a CEO, Chris.
If you were \textbf{the} CEO of -- I guess you still are.
If you were \textbf{the} CEO of an electric utility, and you told your executive team,"I want to build a nuclear power plant,"two of \textbf{the} first questions you would ask are, number one : How much will it cost?
And there ' s not a single engineering consulting firm that I ' ve been able to find anywhere in \textbf{the} world that will put their name on an opinion giving you a cost estimate.
They just don ' t know.
A second question you would ask is : How long will it take to build it, so we can start selling \textbf{the} electricity?
And again, \textbf{the} answer you will get is,"We have no idea."So if you don ' t know how much it ' s going to cost, and you don ' t know when it ' s going to be finished, and you already know that \textbf{the} electricity is more expensive than \textbf{the} alternate ways to produce it, that ' s going to be a little discouraging, and, in fact, that ' s been \textbf{the} case for utilities around \textbf{the} world.
CA : OK.
So there ' s definitely an interesting debate there, but we ' re going to come on to some community questions.
Let ' s have \textbf{the} first of those questions up, please.
From Prosanta Chakrabarty :"People who are skeptical of COVID and of climate change seem to be skeptical of science in general.
It may be that \textbf{the} singular message from scientists gets diluted and convoluted.
How do we fix that?"AG : Yeah, that ' s a great question, Prosanta.
Boy, I ' m trying to put this succinctly and shortly.
I think that there has been a feeling that experts in general have kind of let \textbf{the} US down, and that feeling is much more pronounced in \textbf{the} US than in most other countries.
And I think that \textbf{the} considered opinion of what we call experts has been diluted over \textbf{the} last few decades by \textbf{the} unhealthy dominance of big money in our political system, which has found ways to really twist economic policy to benefit elites.
And this sounds a little radical, but it ' s actually what has happened.
And we have gone for more than 40 years without any meaningful increase in middle-income pay, and where \textbf{the} injustice experienced by African Americans and other communities of color are concerned, \textbf{the} differential in pay between African Americans and majority Americans is \textbf{the} same as it was in 1968, and \textbf{the} family wealth, \textbf{the} net worth -- it takes 11 and a half so-called"typical"African American families to make up \textbf{the} net worth of one so-called"typical"White American family.
And you look at \textbf{the} soaring incomes in \textbf{the} top one or \textbf{the} top one-tenth of one percent, and people say,"Wait a minute.
Whoever \textbf{the} experts were that designed these policies, they haven ' t been doing a good job for me."A final point, Chris : there has been an assault on reason.
There has been a war against truth.
There has been a strategy, maybe it was best known as a strategy decades ago by \textbf{the} tobacco companies who hired actors and dressed them up as doctors to falsely reassure people that there were no health consequences from smoking cigarettes, and a hundred million people died as a result.
That same strategy of diminishing \textbf{the} significance of truth, diminishing, as someone said, \textbf{the} authority of knowledge, I think that has made it kind of open season on any inconvenient truth -- forgive another buzz phrase, but it is apt.
We cannot abandon our devotion to \textbf{the} best available evidence tested in reasoned discourse and used as \textbf{the} basis for \textbf{the} best policies we can form.
CA : Is it possible, Al, that one consequence of \textbf{the} pandemic is actually a growing number of people have revisited their opinions on scientists?
I mean, you ' ve had a chance in \textbf{the} last few months to say,"Do I trust my political leader or do I trust this scientist in terms of what they ' re saying about this virus?"Maybe lessons from that could be carried forward?
AG : Well, you know, I think if \textbf{the} polling is accurate, people do trust their doctors a lot more than some of \textbf{the} politicians who seem to have a vested interest in pretending \textbf{the} pandemic isn ' t real.
And if you look at \textbf{the} incredible bust at President Trump ' s rally in Tulsa, a stadium of 19, 000 people with less than \textbf{one-third} filled, according to \textbf{the} fire marshal, you saw all \textbf{the} empty seats if you saw \textbf{the} news \textbf{clips}, so even \textbf{the} most loyal Trump supporters must have decided to trust their doctors and \textbf{the} medical advice rather than Dr.
Donald Trump.
CA : With a little help from \textbf{the} TikTok generation, perchance.
AG : Well, but that didn ' t affect \textbf{the} turnout.
What they did, very cleverly, and I ' m cheering them on, what they did was affect \textbf{the} Trump White House ' s expectations.
They ' re \textbf{the} reason why he went out a couple days beforehand and said,"We ' ve had a million people sign up."But they didn ' t prevent -- they didn ' t take seats that others could have otherwise taken.
They didn ' t affect \textbf{the} turnout, just \textbf{the} expectations.
CA : OK, let ' s have our next question here."Are you concerned \textbf{the} world will rush back to \textbf{the} use of \textbf{the} private car out of fear of using shared public \textbf{transportation}?"AG : Well, that could actually be one of \textbf{the} consequences, absolutely.
Now, \textbf{the} trends on mass transit were already \textbf{inching} in \textbf{the} wrong direction because of Uber and Lyft and \textbf{the} ridesharing services, and if autonomy ever reaches \textbf{the} goals that its advocates have hoped for then that may also have a similar effect.
But there ' s no doubt that some people are going to be probably a little more reluctant to take mass \textbf{transportation} until \textbf{the} fear of this pandemic is well and truly gone.
CA : Yeah.
Might need a vaccine on that one.
AG : Yeah.
CA : Next question.
Sonaar Luthra, thank you for this question from LA."Given \textbf{the} temperature rise in \textbf{the} Arctic this past week, seems like \textbf{the} rate we are losing our \textbf{carbon} sinks like permafrost or forests is accelerating faster than we predicted.
Are our models too focused on human \textbf{emissions}?"Interesting question.
AG : Well, \textbf{the} models are focused on \textbf{the} factors that have led to these incredible temperature spikes in \textbf{the} north of \textbf{the} Arctic Circle.
They were predicted, they have been predicted, and one of \textbf{the} reasons for it is that as \textbf{the} snow and ice cover melts, \textbf{the} sun ' s incoming rays are no longer reflected back into space at a 90 percent rate, and instead, when they fall on \textbf{the} dark tundra or \textbf{the} dark ocean, they ' re absorbed at a 90 percent rate.
So that ' s a magnifier of \textbf{the} warming in \textbf{the} Arctic, and this has been predicted.
There are a number of other consequences that are also in \textbf{the} models, but some of them may have to be recalibrated. \textbf{The} scientists are freshly concerned that \textbf{the} \textbf{emissions} of both CO2 and methane from \textbf{the} thawing tundra could be larger than they had hoped they would be.
There ' s also just been a brand-new study.
I won ' t spend time on this, because it deals with a kind of geeky term called"climate sensitivity,"which has been a factor in \textbf{the} models with large error bars because it ' s so hard to pin down.
But \textbf{the} latest evidence indicates, worryingly, that \textbf{the} sensitivity may be greater than they had thought, and we will have an even more daunting task.
That shouldn ' t discourage us.
I truly believe that once we cross this tipping point, and I do believe we ' re doing it now, as I ' ve said, then I think we ' re going to find a lot of ways to speed up \textbf{the} \textbf{emissions} reductions.
CA : We ' ll take one more question from \textbf{the} community.
Haha."Geoengineering is making extraordinary progress.
Exxon is investing in technology from Global Thermostat that seems promising.
What do you think of these air and water \textbf{carbon} capture technologies?"Stephen Petranek.
AG : Yeah.
Well, you and I have talked about this before, Chris.
I ' ve been strongly opposed to conducting an unplanned global experiment that could go wildly wrong, and most are really scared of that approach.
However, \textbf{the} term"geoengineering"is a nuanced term that covers a lot.
If you want to paint roofs white to reflect more energy from \textbf{the} cityscapes, that ' s not going to bring a danger of a runaway effect, and there are some other things that are loosely called"geoengineering"like that, which are fine.
But \textbf{the} idea of blocking out \textbf{the} sun ' s rays -- that ' s insane in my opinion.
Turns out plants need sunlight for photosynthesis and solar panels need sunlight for producing electricity from \textbf{the} sun ' s rays.
And \textbf{the} consequences of changing everything we know and pretending that \textbf{the} consequences are going to precisely cancel out \textbf{the} unplanned experiment of global warming that we already have underway, you know, there are glitches in our thinking.
One of them is called \textbf{the}"single solution bias,"and there are people who just have a hunger to say,"Well, that one solution, we just need to latch on to that and do that, and damn \textbf{the} consequences."Well, it ' s \textbf{nuts}.
CA : But let me push back on this just a little bit.
So let ' s say that we agree that a single solution, all-or-nothing attempt at geoengineering is crazy.
But there are scenarios where \textbf{the} world looks at \textbf{emissions} and just sees, in 10 years ' time, let ' s say, that they are just not coming down fast enough and that we are at risk of several other liftoff events where this train will just get away from us, and we will see temperature rises of three, four, five, six, seven degrees, and all of civilization is at risk.
Surely, there is an approach to geoengineering that could be modeled, in a way, on \textbf{the} way that we approach medicine.
Like, for hundreds of years, we don ' t really understand \textbf{the} human body, people would try interventions, and some of them would work, and some of them wouldn ' t.
No one says in medicine,"You know, go in and take an all-or-nothing decision on someone ' s life,"but they do say,"Let ' s try some stuff."If an experiment can be reversible, if it ' s plausible in \textbf{the} first place, if there ' s reason to think that it might work, we actually owe it to \textbf{the} future health of humanity to try at least some types of tests to see what could work.
So, small tests to see whether, for example, seeding of something in \textbf{the} ocean might create, in a nonthreatening way, \textbf{carbon} sinks.
Or maybe, rather than filling \textbf{the} atmosphere with sulfur dioxide, a smaller experiment that was not that big a deal to see whether, cost-effectively, you could reduce \textbf{the} temperature a little bit.
Surely, that isn ' t completely crazy and is at least something we should be thinking about in case these other measures don ' t work?
AG : Well, there ' ve already been such experiments to seed \textbf{the} ocean to see if that can increase \textbf{the} uptake of CO2.
And \textbf{the} experiments were an unmitigated failure, as many predicted they would be.
But that, again, is \textbf{the} kind of approach that ' s very different from putting tinfoil strips in \textbf{the} atmosphere orbiting \textbf{the} Earth.
That was \textbf{the} way that solar geoengineering proposal started.
Now they ' re focusing on chalk, so we have chalk dust all over everything.
But more serious than that is \textbf{the} fact that it might not be reversible.
CA : But, Al, that ' s \textbf{the} rhetoric response. \textbf{The} amount of dust that you need to drop by a degree or two wouldn ' t result in chalk dust over everything.
It would be unbelievably -- like, it would be less than \textbf{the} dust that people experience every day, anyway.
I mean, I just -- AG : First of all, I don ' t know how you do a small experiment in \textbf{the} atmosphere.
And secondly, if we were to take that approach, we would have to steadily increase \textbf{the} amount of whatever substance they decided.
We ' d have to increase it every single year, and if we ever stopped, then there would be a sudden snapback, like"\textbf{The} Picture of Dorian Gray,"that old book and movie, where suddenly all of \textbf{the} things caught up with you at once. \textbf{The} fact that anyone is even considering these approaches, Chris, is a measure of a feeling of desperation that some have begun to feel, which I understand, but I don ' t think it should drive us toward these reckless experiments.
And by \textbf{the} way, using your analogy to experimental cancer treatments, for example, you usually get informed consent from \textbf{the} patient.
Getting informed consent from 7. 8 billion people who have no voice and no say, who are subject to \textbf{the} potentially catastrophic consequences of this wackadoodle proposal that somebody comes up with to try to rearrange \textbf{the} entire Earth ' s atmosphere and hope and pretend that it ' s going to cancel out, \textbf{the} fact that we ' re putting 152 million tons of heat-trapping, manmade global warming pollution into \textbf{the} sky every day.
That ' s what ' s really insane.
A scientist decades ago compared it this way.
He said, if you had two people on a sinking boat and one of them says,"You know, we could probably use some mirrors to signal to shore to get them to build a sophisticated wave-generating machine that will cancel out \textbf{the} \textbf{rocking} of \textbf{the} boat by these guys in \textbf{the} back of \textbf{the} boat."Or you could get them to stop \textbf{rocking} \textbf{the} boat.
And that ' s what we need to do.
We need to stop what ' s causing \textbf{the} crisis.
CA : Yeah, that ' s a great story, but if \textbf{the} effort to stop \textbf{the} people \textbf{rocking} in \textbf{the} back of \textbf{the} boat is as complex as \textbf{the} scientific proposal you just outlined, whereas \textbf{the} experiment to stop \textbf{the} waves is actually as simple as telling \textbf{the} people to stop \textbf{rocking} \textbf{the} boat, that story changes.
And I think you ' re right that \textbf{the} issue of informed consent is a really challenging one, but, I mean, no one gave informed consent to do all of \textbf{the} other things we ' re doing to \textbf{the} atmosphere.
And I agree that \textbf{the} moral hazard issue is worrying, that if we became dependent on geoengineering and took away our efforts to do \textbf{the} rest, that would be tragic.
It just seems like, I wish it was possible to have a nuanced debate of people saying, you know what, there ' s multiple dials to a very complex problem.
We ' re going to have to adjust several of them very, very carefully and keep talking to each other.
Wouldn ' t that be a goal to just try and have a more nuanced debate about this, rather than all of that geoengineering can ' t work?
AG : Well, I ' ve said some of it, you know, \textbf{the} benign forms that I ' ve mentioned, I ' m not ruling those out.
But blocking \textbf{the} Sun ' s rays from \textbf{the} Earth, not only do you affect 7. 8 billion people, you affect \textbf{the} plants and \textbf{the} animals and \textbf{the} ocean currents and \textbf{the} wind currents and natural processes that we ' re in danger of disrupting even more.
Techno-optimism is something I ' ve engaged in in \textbf{the} past, but to latch on to some brand-new technological solution to rework \textbf{the} entire Earth ' s natural system because somebody thinks he ' s clever enough to do it in a way that precisely cancels out \textbf{the} consequences of using \textbf{the} atmosphere as an open sewer for heat-trapping manmade gases.
It ' s much more important to stop using \textbf{the} atmosphere as an open sewer.
That ' s what \textbf{the} problem is.
CA : All right, well, we ' ll agree that that is \textbf{the} most important thing, for sure, and speaking of which, do you believe \textbf{the} world needs \textbf{carbon} pricing, and is there any prospect for getting there?
AG : Yes.
Yes to both questions.
For decades, almost every economist who is asked about \textbf{the} climate crisis says,"Well, we just need to put a price on \textbf{carbon}."And I have certainly been in favor of that approach.
But it is daunting.
Nevertheless, there are 43 jurisdictions around \textbf{the} world that already have a price on \textbf{carbon}.
We ' re seeing it in Europe.
They finally straightened out their \textbf{carbon} pricing mechanism.
It ' s an \textbf{emissions} trading version of it.
We have places that have put a tax on \textbf{carbon}.
That ' s \textbf{the} approach \textbf{the} economists prefer.
China is beginning to implement its national \textbf{emissions} trading program.
California and quite a few other states in \textbf{the} US are already doing it.
It can be given back to people in a revenue-neutral way.
But \textbf{the} opposition to it, Chris, which you ' ve noted, is impressive enough that we do have to take other approaches, and I would say most climate activists are now saying, look, let ' s don ' t make \textbf{the} best \textbf{the} enemy of \textbf{the} better.
There are other ways to do this as well.
We need every solution we can rationally employ, including by regulation.
And often, when \textbf{the} political difficulty of a proposal becomes too difficult in a market-oriented approach, \textbf{the} fallback is with regulation, and it ' s been given a bad name, regulation, but many places are doing it.
I mentioned phasing out internal combustion engines.
That ' s an example.
There are 160 cities in \textbf{the} US that have already by regulation ordered that within a date certain, 100 percent of all their electricity will have to come from \textbf{renewable} sources.
And again, \textbf{the} market forces that are driving \textbf{the} cost of \textbf{renewable} energy and \textbf{sustainability} solutions ever downward, that gives us \textbf{the} wind at our back.
This is working in our favor.
CA : I mean, \textbf{the} pushback on \textbf{carbon} pricing often goes further from parts of \textbf{the} environmental movement, which is to a pushback on \textbf{the} role of business in general.
Business is actually -- well, capitalism -- is blamed for \textbf{the} climate crisis because of unrelenting growth, to \textbf{the} point where many people don ' t trust business to be part of \textbf{the} solution. \textbf{The} only way to go forward is to regulate, to force businesses to do \textbf{the} right thing.
Do you think that business has to be part of \textbf{the} solution?
AG : Well, definitely, because \textbf{the} allocation of capital needed to solve this crisis is greater than what governments can handle.
And businesses are beginning, many businesses are beginning to play a very constructive role.
They ' re getting a demand that they do so from their customers, from their investors, from their boards, from their executive teams, from their families.
And by \textbf{the} way, \textbf{the} rising generation is demanding a brighter future, and when CEOs interview potential new hires, they find that \textbf{the} new hires are interviewing them.
They want to make a nice income, but they want to be able to tell their family and friends and peers that they ' re doing something more than just making money.
One illustration of how this new generation is changing, Chris : there are 65 colleges in \textbf{the} US right now where \textbf{the} College Young Republican Clubs have joined together to jointly demand that \textbf{the} Republican National Committee change its policy on climate, lest they lose that entire generation.
This is a global phenomenon. \textbf{The} Greta Generation is now leading this in so many ways, and if you look at \textbf{the} polling, again, \textbf{the} vast majority of young Republicans are demanding a change on climate policy.
This is really a movement that is building still.
CA : I was going to ask you about that, because one of \textbf{the} most painful things over \textbf{the} last 20 years has just been how climate has been politicized, certainly in \textbf{the} US.
You ' ve probably felt yourself at \textbf{the} heart of that a lot of \textbf{the} time, with people attacking you personally in \textbf{the} most merciless, and unfair ways, often.
Do you really see signs that that might be changing, led by \textbf{the} next generation?
AG : Yeah, there ' s no question about it.
I don ' t want to rely on polls too much.
I ' ve mentioned them already.
But there was a new one that came out that looked at \textbf{the} wavering Trump supporters, those who supported him strongly in \textbf{the} past and want to do so again. \textbf{The} number one issue, surprisingly to some, that is giving them pause, is \textbf{the} craziness of President Trump and his administration on climate.
We ' re seeing big majorities of \textbf{the} Republican Party overall saying that they ' re ready to start exploring some real solutions to \textbf{the} climate crisis.
I think that we ' re really getting there, no question about it.
CA : I mean, you ' ve been \textbf{the} figurehead for raising this issue, and you happen to be a Democrat.
Is there anything that you can personally do to -- I don ' t know -- to open \textbf{the} tent, to welcome people, to try and say,"This is beyond politics, dear friends"?
AG : Yeah.
Well, I ' ve tried all of those things, and maybe it ' s made a little positive difference.
I ' ve worked with \textbf{the} Republicans extensively.
And, you know, well after I left \textbf{the} White House, I had Newt Gingrich and Pat Robertson and other prominent Republicans appear on national TV ads with me saying we ' ve got to solve \textbf{the} climate crisis.
But \textbf{the} petroleum industry has really doubled down enforcing discipline within \textbf{the} Republican Party.
I mean, look at \textbf{the} attacks they ' ve launched against \textbf{the} Pope when he came out with his encyclical and was demonized, not by all for sure, but there were hawks in \textbf{the} anti-climate movement who immediately started training their guns on Pope Francis, and there are many other examples.
They enforce discipline and try to make it a partisan issue, even as Democrats reach out to try to make it bipartisan.
I totally agree with you that it should not be a partisan issue.
It didn ' t use to be, but it ' s been artificially weaponized as an issue.
CA : I mean, \textbf{the} CEOs of \textbf{oil} companies also have kids who are talking to them.
It feels like some of them are moving and are trying to invest and trying to find ways of being part of \textbf{the} future.
Do you see signs of that?
AG : Yeah.
I think that business leaders, including in \textbf{the} \textbf{oil} and gas companies, are hearing from their families.
They ' re hearing from their friends.
They ' re hearing from their employees.
And, by \textbf{the} way, we ' ve seen in \textbf{the} tech industry some mass walkouts by employees who are demanding that some of \textbf{the} tech companies do more and get serious.
I ' m so proud of Apple.
Forgive me for parenthetically praising Apple.
You know, I ' m on \textbf{the} board, but I ' m such a big fan of Tim Cook and my colleagues at Apple.
It ' s an example of a tech company that ' s really doing \textbf{fantastic} things.
And there ' s some others as well.
There are others in many industries.
But \textbf{the} pressures on \textbf{the} \textbf{oil} and gas companies are quite extraordinary.
You know, BP just wrote down 12 and a half billion dollars ' worth of \textbf{oil} and gas assets and said that they ' re never going to see \textbf{the} light of day. \textbf{Two-thirds} of \textbf{the} \textbf{fossil} fuels that have already been discovered cannot be burned and will not be burned.
And so that ' s a big economic risk to \textbf{the} global economy, like \textbf{the} subprime mortgage crisis.
We ' ve got 22 trillion dollars of subprime \textbf{carbon} assets, and just yesterday, there was a major report that \textbf{the} fracking industry in \textbf{the} US is seeing now a wave of bankruptcies because \textbf{the} price of \textbf{the} fracked gas and \textbf{oil} has fallen below levels that make them economic.
CA : Is \textbf{the} shorthand of what ' s happened there that electric cars and electric technologies and solar and so forth have helped drive down \textbf{the} price of \textbf{oil} to \textbf{the} point where huge amounts of \textbf{the} reserves just can ' t be developed profitably?
AG : Yes, that ' s it.
That ' s mainly it. \textbf{The} projections for energy sources in \textbf{the} next several years uniformly predict that electricity from wind and solar is going to continue to plummet in price, and therefore using gas or \textbf{coal} to make steam to turn \textbf{the} \textbf{turbines} is just not going to be economical.
Similarly, \textbf{the} electrification of \textbf{the} \textbf{transportation} sector is having \textbf{the} same effect.
Some are also looking at \textbf{the} trend in national, regional and local governance.
I mentioned this before, but they ' re predicting a very different energy future.
But let me come back, Chris, because we talked about business leaders.
I think you were getting in a question a moment ago about capitalism itself, and I do want to say a word on that, because there are a lot of people who say maybe capitalism is \textbf{the} basic problem.
I think \textbf{the} current form of capitalism we have is desperately in need of reform. \textbf{The} short-term outlook is often mentioned, but \textbf{the} way we measure what is of value to us is also at \textbf{the} heart of \textbf{the} crisis of modern capitalism.
Now, capitalism is at \textbf{the} base of every successful economy, and it balances supply and demand, \textbf{unlocks} a higher \textbf{fraction} of \textbf{the} human potential, and it ' s not going anywhere, but it needs to be reformed, because \textbf{the} way we measure what ' s valuable now ignores so-called negative externalities like pollution.
It also ignores positive externalities like investments in education and health care, mental health care, family services.
It ignores \textbf{the} depletion of resources like groundwater and topsoil and \textbf{the} \textbf{web} of living species.
And it ignores \textbf{the} distribution of incomes and net worths, so when GDP goes up, people cheer, two percent, three percent -- wow! -- four percent, and they think,"Great!"But it ' s accompanied by vast increases in pollution, chronic underinvestment in public goods, \textbf{the} depletion of irreplaceable natural resources, and \textbf{the} worst inequality crisis we ' ve seen in more than a hundred years that is threatening \textbf{the} future of both capitalism and democracy.
So we have to change it.
We have to reform it.
CA : So reform capitalism, but don ' t throw it out.
We ' re going to need it as a tool as we go forward if we ' re to solve this.
AG : Yeah, I think that ' s right, and just one other point : \textbf{the} worst environmental abuses in \textbf{the} last hundred years have been in jurisdictions that experimented during \textbf{the} 20th century with \textbf{the} alternatives to capitalism on \textbf{the} \textbf{left} and right.
CA : Interesting.
All right.
Two last community questions quickly.
Chadburn Blomquist :"As you are reading \textbf{the} tea leaves of \textbf{the} impact of \textbf{the} current pandemic, what do you think in regard to our response to combatting climate change will be \textbf{the} most impactful lesson learned?"AG : Boy, that ' s a very thoughtful question, and I wish my answer could rise to \textbf{the} same level on short notice.
I would say first, don ' t ignore \textbf{the} scientists.
When there is virtual unanimity among \textbf{the} scientific and medical experts, pay attention.
Don ' t let some politician dissuade you.
I think President Trump is slowly learning that ' s it ' s kind of difficult to gaslight a virus.
He tried to gaslight \textbf{the} virus in Tulsa.
It didn ' t come off very well, and tragically, he decided to recklessly roll \textbf{the} dice a month ago and ignore \textbf{the} recommendations for people to wear masks and to socially distance and to do \textbf{the} other things, and I think that lesson is beginning to take hold in a much stronger way.
But beyond that, Chris, I think that this period of time has been characterized by one of \textbf{the} most profound opportunities for people to rethink \textbf{the} patterns of their lives and to consider whether or not we can ' t do a lot of things better and differently.
And I think that this rising generation I mentioned before has been even more profoundly affected by this interlude, which I hope ends soon, but I hope \textbf{the} lessons endure.
I expect they will.
CA : Yeah, it ' s amazing how many things you can do without emitting \textbf{carbon}, that we ' ve been forced to do.
Let ' s have one more question here.
Frank Hennessy :"Are you encouraged by \textbf{the} ability of people to quickly adapt to \textbf{the} new normal due to COVID-19 as evidence that people can and will change their habits to respond to climate change?"AG : Yes, but I think we have to keep in mind that there is a crisis within this crisis. \textbf{The} impact on \textbf{the} African American community, which I mentioned before, on \textbf{the} Latinx community, Indigenous peoples. \textbf{The} highest infection rate is in \textbf{the} Navajo Nation right now.
So some of these questions appear differently to those who are really getting \textbf{the} brunt of this crisis, and it is unacceptable that we allow this to continue.
It feels one way to you and me and perhaps to many in our audience today, but for low-income communities of color, it ' s an entirely different crisis, and we owe it to them and to all of us to get busy and to start using \textbf{the} best science and solve this pandemic.
You know \textbf{the} phrase"pandemic economics."Somebody said, \textbf{the} first principle of pandemic economics is take care of \textbf{the} pandemic, and we ' re not doing that yet.
We ' re seeing \textbf{the} president try to goose \textbf{the} economy for his reelection, never mind \textbf{the} prediction of tens of thousands of additional American deaths, and that is just unforgivable in my opinion.
CA : Thank you, Frank.
So Al, you, along with others in \textbf{the} community played a key role in encouraging TED to launch this initiative called"Countdown."Thank you for that, and I guess this conversation is continuing among many of us.
If you ' re interested in climate, watching this, check out \textbf{the} Countdown website, countdown. ted. com, and be part of 10 / 10 / 2020, when we are trying to put out an alert to \textbf{the} world that climate can ' t wait, that it really matters, and there ' s going to be some amazing content free to \textbf{the} world on that day.
Thank you, Al, for your inspiration and support in doing that.
I wonder whether you could end today ' s session just by painting us a picture, like how might things roll out over \textbf{the} next decade or so?
Just tell us whether there is still a story of hope here.
AG : I ' d be glad to.
I ' ve got to get one plug in.
I ' ll make it brief.
July 18 through July 26, \textbf{The} Climate Reality Project is having a global training.
We ' ve already had 8, 000 people register.
You can go to climatereality. com.
Now, a bright future.
It begins with all of \textbf{the} kinds of efforts that you ' ve thrown yourself into in organizing Countdown.
Chris, you and your team have been amazing to work with, and I ' m so excited about \textbf{the} Countdown project.
TED has an unparalleled ability to spread ideas that are worth spreading, to raise consciousness, to enlighten people around \textbf{the} world, and it ' s needed for climate and \textbf{the} solutions to \textbf{the} climate crisis like it ' s never been needed before, and I just want to thank you for what you personally are doing to organize this \textbf{fantastic} Countdown program.
CA : Thank you.
And \textbf{the} world?
Are we going to do this?
Do you think that humanity is going to pull this off and that our grandchildren are going to have beautiful lives where they can celebrate nature and not spend every day in fear of \textbf{the} next tornado or \textbf{tsunami}?
AG : I am optimistic that we will do it, but \textbf{the} answer is in our hands.
We have seen dark times in periods of \textbf{the} past, and we have risen to meet \textbf{the} challenge.
We have limitations of our long \textbf{evolutionary} heritage and elements of our culture, but we also have \textbf{the} ability to transcend our limitations, and when \textbf{the} chips are down, and when survival is at stake and when our children and future generations are at stake, we ' re capable of more than we sometimes allow ourselves to think we can do.
This is such a time.
I believe we will rise to \textbf{the} occasion, and we will create a bright, clean, prosperous, just and fair future.
I believe it with all my heart.
CA : Al Gore, thank you for your life of work, for all you ' ve done to elevate this issue and for spending this time with us now.
Thank you.
AG : Back at you.
Thank you.

% matched lemmas: agriculture, carbon, clip, coal, emission, evolutionary, fantastic, fossil, fraction, greenhouse, inch, left, nut, occur, oil, one-third, planet, renewable, rock, sustainability, technique, the, transportation, tsunami, turbine, two-third, unlock, web
\end{textsample}
